\documentclass[12pt]{article}
\usepackage[utf8]{inputenc}
\usepackage[T1]{fontenc}
\usepackage[margin=1in]{geometry}
\usepackage{setspace}
\usepackage{xcolor}
\usepackage[colorlinks=true, urlcolor=blue!60!black, linkcolor=blue!60!black]{hyperref}
\setlength{\parindent}{2em}
\setlength{\parskip}{1.25em}
\renewcommand{\baselinestretch}{1.0}
\begin{document}
\begin{center}
{\LARGE \textit{Pierce v. Society of Sisters}, 268 U.S. 510 (1925)\par}
\vspace{0.8em}
{\large Case Brief\par}
\end{center}
\vspace{1em}

\section{Memory Jogger}
A unanimous Supreme Court struck down Oregon's compulsory public-school law under the Fourteenth Amendment's Due Process Clause, holding that the state cannot ``standardize its children'' by abolishing private and parochial schools because parents hold a fundamental liberty interest to direct their children's upbringing and education.

\section{Detailed Case Facts}
On November 7, 1922, amid post-World War I nativism and widespread anti-Catholic sentiment backed by groups including the Ku Klux Klan and the Federation of Patriotic Societies, Oregon voters adopted the Compulsory Education Act by ballot initiative, amending Section 5259 of the Oregon Laws (\href{https://supreme.justia.com/cases/federal/us/268/510/}{Justia}; \href{https://www.oregonencyclopedia.org/packets/1}{Oregon Encyclopedia}). Scheduled to take effect on September 1, 1926, the statute required every parent, guardian, or other person having control or custody of a child between the ages of eight and sixteen to send the child ``to a public school for the period of time a public school shall be held during the current year'' in the child's resident district. Violation was a misdemeanor punishable by a fine of \$5 to \$100, imprisonment in the county jail for two to thirty days, or both, with each day's failure constituting a separate offense (\href{https://supreme.justia.com/cases/federal/us/268/510/}{Justia}). Narrow statutory exemptions applied only to children who were physically or mentally unable to attend, had completed the eighth grade, lived beyond specified road distances from a public school (unless transportation was provided), or received private tutoring with annual written permission from, and quarterly examinations by, the county superintendent. The law eliminated the prior exemption for state-recognized private schools, effectively outlawing private and parochial primary education across Oregon (\href{https://www.courtlistener.com/opinion/100694/pierce-v-society-of-sisters/}{CourtListener}; \href{https://supreme.justia.com/cases/federal/us/268/510/}{Justia}).

Two Oregon corporations operating private schools filed separate suits to enjoin enforcement of the Act against Oregon Governor Walter M. Pierce, Attorney General Isaac H. Van Winkle, and Multnomah County District Attorney Stanley Myers (\href{https://www.oyez.org/cases/1900-1940/268us510}{Oyez}; \href{https://supreme.justia.com/cases/federal/us/268/510/}{Justia}). First, the Society of the Sisters of the Holy Names of Jesus and Mary (incorporated in 1880) operated interdependent Roman Catholic primary schools, high schools, junior colleges, and orphanages in Oregon, teaching standard Oregon public-school subjects alongside systematic Catholic religious instruction and moral training. Its primary schools generated more than \$30,000 in annual income and depended on long-term contracts with teachers and parents. Second, Hill Military Academy (incorporated in 1908) was a for-profit secular elementary, college-preparatory, and military training school for boys aged five to twenty-one, averaging 100 students at roughly \$800 in annual fees per student, with a curriculum conforming to State Board of Education standards (\href{https://supreme.justia.com/cases/federal/us/268/510/}{Justia}). Even before the Act's 1926 effective date, state officials' public declarations that they would strictly enforce the law caused parents to withdraw students and refuse future enrollment contracts, steadily eroding both schools' enrollments, income, and property values.

\section{Procedural History}
The Society of Sisters (No.~583) and Hill Military Academy (No.~584) filed bills in equity in the U.S. District Court for the District of Oregon seeking injunctive relief on the ground that the Act violated the Fourteenth Amendment's Due Process Clause. The defendants interposed no answers contesting the factual allegations. Pursuant to Section 266 of the Judicial Code, a three-judge federal district court heard the motions for preliminary injunctions and ruled for both schools, holding that the Fourteenth Amendment protected their property and business goodwill against unlawful state interference with the free choice of parents and guardians, and that the suits were not premature because irreparable injury was already occurring (\textit{Society of the Sisters of the Holy Names of Jesus and Mary v. Pierce}, 296 F. 928 (D. Or. 1924)) (\href{https://supreme.justia.com/cases/federal/us/268/510/}{Justia}). Oregon officials appealed directly to the U.S. Supreme Court. The Court heard consolidated oral argument on March 16--17, 1925, and unanimously affirmed the District Court's decrees on June 1, 1925 (\href{https://www.oyez.org/cases/1900-1940/268us510}{Oyez}; \href{https://supreme.justia.com/cases/federal/us/268/510/}{Justia}).

\section{Judicial Votes}
The decision was \textbf{unanimous (9--0)} (\href{https://constitutioncenter.org/the-constitution/supreme-court-case-library/pierce-v-society-of-sisters}{National Constitution Center}; \href{https://www.oyez.org/cases/1900-1940/268us510}{Oyez}):
\begin{itemize}
    \item \textbf{Opinion of the Court:} Justice James Clark McReynolds, joined by Chief Justice William Howard Taft and Associate Justices Oliver Wendell Holmes, Jr., Willis Van Devanter, Louis D. Brandeis, George Sutherland, Pierce Butler, Edward T. Sanford, and Harlan Fiske Stone.
    \item \textbf{Concurrences or Dissents:} None.
\end{itemize}

\section{Holding}
\textbf{Judgment:} Affirmed. The preliminary injunctions barring Oregon officials from enforcing the Compulsory Education Act of 1922 were upheld (\href{https://supreme.justia.com/cases/federal/us/268/510/}{Justia}).

\textbf{Rule Announced:} Under the Due Process Clause of the Fourteenth Amendment and the doctrine of \textit{Meyer v. Nebraska}, 262 U.S. 390 (1923), a state may not compel all normal children to attend public schools exclusively, because doing so unreasonably interferes with the fundamental liberty of parents and guardians to direct the upbringing and education of children under their control without any reasonable relation to a purpose within the competency of the state (\href{https://supreme.justia.com/cases/federal/us/268/510/}{Justia}). Furthermore, although corporations cannot claim Fourteenth Amendment ``liberty'' for themselves, private school corporations have direct, immediate standing to seek equitable relief before a statute's effective date when unconstitutional state compulsion over their patrons threatens immediate, irreparable destruction of their lawful business and property (\href{https://supreme.justia.com/cases/federal/us/268/510/}{Justia}).

\textbf{Holding vs.\ Dicta:} The \textit{holding} invalidates Oregon's total ban on private primary school attendance as an unconstitutional deprivation of parental liberty and corporate property without due process of law, and confirms both pre-enforcement ripeness and the schools' standing to assert the injury flowing from unlawful compulsion of parents. The Court's famous concession that ``[n]o question is raised concerning the power of the State reasonably to regulate all schools, to inspect, supervise and examine them, their teachers and pupils; to require that all children of proper age attend some school, that teachers shall be of good moral character and patriotic disposition, that certain studies plainly essential to good citizenship must be taught, and that nothing be taught which is manifestly inimical to the public welfare'' (268 U.S. at 534) is technically \textit{dicta}, because Oregon had not enacted a regulatory scheme for private schools---it had abolished them. Nevertheless, this dicta became foundational in later cases defining the boundary between parental rights and state educational regulation, including \textit{Wisconsin v. Yoder}, 406 U.S. 205 (1972) (\href{https://supreme.justia.com/cases/federal/us/268/510/}{Justia}; \href{https://fedsoc.org/case/pierce-v-society-of-sisters}{Federalist Society}).

\section{Analysis of Opinions}

\subsection{Opinion of the Court (Justice McReynolds)}
Justice McReynolds wrote a concise, forceful opinion for a unanimous Court that blended \textit{Lochner}-era economic property protection with an expansive vision of non-economic personal liberty (\href{https://supreme.justia.com/cases/federal/us/268/510/}{Justia}; \href{https://www.courtlistener.com/opinion/100694/pierce-v-society-of-sisters/}{CourtListener}). He began by bracketing what was \textit{not} at issue: the state's unquestioned police power to require compulsory education at \textit{some} school and to reasonably inspect, supervise, and set basic curricular and teacher-character standards for all schools. Turning to the practical impact of Oregon's 1922 Act, he observed that enforcement would inevitably destroy the appellees' primary schools---and likely all private primary schools for normal children in Oregon---even though the schools were engaged in a ``useful and meritorious'' enterprise and no emergency justified extraordinary state intervention (\href{https://supreme.justia.com/cases/federal/us/268/510/}{Justia}).

Relying directly on \textit{Meyer v. Nebraska}, 262 U.S. 390 (1923), Justice McReynolds held that the statute ``unreasonably interferes with the liberty of parents and guardians to direct the upbringing and education of children under their control.'' Rights guaranteed by the Constitution, he explained, may not be abridged by legislation bearing no reasonable relation to a purpose within the state's competency. In the opinion's most celebrated passage, he declared that ``[t]he fundamental theory of liberty upon which all governments in this Union repose excludes any general power of the State to standardize its children by forcing them to accept instruction from public teachers only,'' emphasizing that ``[t]he child is not the mere creature of the State'' (\href{https://supreme.justia.com/cases/federal/us/268/510/}{Justia}).

Justice McReynolds then addressed three doctrinal objections raised by Oregon. First, while conceding that corporations cannot claim Fourteenth Amendment ``liberty'' for themselves, he held that the Society of Sisters and Hill Military Academy possessed ``business and property'' threatened with destruction by the state's unconstitutional compulsion of their patrons, entitling them to equitable relief under precedents such as \textit{Truax v. Raich}, 239 U.S. 33 (1915), and \textit{Terrace v. Thompson}, 263 U.S. 197 (1923) (\href{https://supreme.justia.com/cases/federal/us/268/510/}{Justia}). Second, because the statute applied generally to all parents and schools, it could not be defended as an exercise of the state's reserved power to amend corporate charters (distinguishing \textit{Berea College v. Kentucky}, 211 U.S. 45 (1908)). Third, the suits were not premature even though the Act did not take effect until September 1926, because the threat of enforcement was already causing present, real, and irreparable loss of students and revenue.

\section{Examples: Future Applications}
\begin{enumerate}
    \item \textbf{Same-Side (Unconstitutional):} To promote civic cohesion, a state enacts a law prohibiting all secular and religious private middle schools and requiring every child aged eleven to fourteen to attend a state-run public middle school, imposing misdemeanor fines on noncompliant parents. \textit{Analysis:} This statute directly replicates Oregon's Compulsory Education Act by attempting to ``standardize'' children through a state monopoly on schooling, violating parents' Fourteenth Amendment substantive due process liberty under \textit{Pierce}.
    \item \textbf{Same-Side (Unconstitutional):} A state legislature bars parents from enrolling their children in accredited private boarding academies outside the state unless the local public school superintendent certifies that the local public school cannot meet the child's academic needs. \textit{Analysis:} By conditioning access to lawful private education on state permission and effectively forcing families into public schools, this law unreasonably interferes with parental control over education without a valid state-competency justification under \textit{Pierce}.
    \item \textbf{Opposite-Side (Constitutional):} A state requires all private and parochial elementary schools to comply with municipal fire-safety codes, maintain basic immunization records, operate at least 180 instructional days per year, and teach core subjects of reading, mathematics, science, and U.S. civics, while leaving schools free to add religious or specialized instruction. \textit{Analysis:} These neutral health, safety, and baseline curricular rules fall squarely within \textit{Pierce}'s express recognition of the state's power ``reasonably to regulate all schools'' and require ``certain studies plainly essential to good citizenship.''
    \item \textbf{Opposite-Side (Constitutional):} A state excludes private religious schools from receiving direct taxpayer-funded operational grants while continuing to permit parents to send their children to those schools at private expense (assuming no discriminatory exclusion from a generally available public benefit program). \textit{Analysis:} \textit{Pierce} establishes a negative liberty right to choose private education free from state prohibition, not an affirmative substantive due process right to compel the state to subsidize private schooling.
    \item \textbf{Fence-Sitter (Uncertain / Close Case):} A state imposes a comprehensive ``substantial equivalency'' regime requiring every private religious school to devote at least five hours per day to secular subjects taught in English by state-certified teachers using state-approved textbooks, effectively threatening to close ultra-Orthodox yeshivas whose curriculum focuses primarily on Talmudic study in Yiddish and Hebrew. \textit{Analysis:} This scenario sits directly on the contested boundary between \textit{Pierce}'s parental-liberty holding (reinforced by Free Exercise claims after \textit{Wisconsin v. Yoder}) and \textit{Pierce}'s regulatory dicta permitting states to enforce baseline citizenship-education standards; courts and scholars remain sharply divided over how intrusive ``substantial equivalency'' mandates may be before they functionally destroy private religious schooling (\href{https://digital.sandiego.edu/cgi/viewcontent.cgi?article=1073&context=jcli}{Journal of Contemporary Legal Issues}).
\end{enumerate}

\section{Critique}
\textbf{Originalist and Textualist Perspectives:} Originalists and textualists have long viewed \textit{Pierce} with methodological ambivalence. Although widely admired in outcome, \textit{Pierce}---alongside \textit{Meyer v. Nebraska}---was born during the \textit{Lochner} era of economic substantive due process and relies on an unenumerated ``liberty'' interest nowhere explicitly mentioned in the Fourteenth Amendment's text. In his dissent in \textit{Troxel v. Granville}, 530 U.S. 57, 91--93 (2000), Justice Scalia acknowledged that parental authority is among the unalienable rights recognized by the Declaration of Independence and retained under the Ninth Amendment, yet he criticized \textit{Meyer} and \textit{Pierce} as relics of ``an era rich in substantive due process holdings that have since been repudiated'' and refused to extend their unenumerated-rights theory, warning that judicially crafted parental rights invite federal courts to displace state legislatures in family law (\href{https://www.law.cornell.edu/supct/html/99-138.ZD1.html}{Scalia Dissent in Troxel, Cornell LII}). Other conservative jurists and scholars---including Justice Kennedy in \textit{Troxel} (530 U.S. at 95)---have argued that had \textit{Pierce} been decided after First Amendment incorporation (which began only one week later in \textit{Gitlow v. New York}, 268 U.S. 652 (1925)), the Court could have grounded the decision more securely in the enumerated Free Exercise and Free Speech Clauses rather than substantive due process (\href{https://firstamendment.mtsu.edu/article/pierce-v-society-of-sisters/}{First Amendment Encyclopedia}; \href{https://fedsoc.org/case/pierce-v-society-of-sisters}{Federalist Society}).

\textbf{Progressive and Doctrinal Perspectives:} Progressive and doctrinal scholars celebrate \textit{Pierce} as a cornerstone of American pluralism and the doctrinal seedbed for modern privacy and autonomy jurisprudence---cited across \textit{Griswold v. Connecticut}, \textit{Roe v. Wade}, \textit{Lawrence v. Texas}, and \textit{Obergefell v. Hodges}---precisely because it rescued personal liberty from the repudiation of \textit{Lochner}'s economic due process (\href{https://constitutioncenter.org/the-constitution/supreme-court-case-library/pierce-v-society-of-sisters}{National Constitution Center}; \href{https://www.educationnext.org/centennial-of-pierce-v-society-of-sisters-landmark-decision-school-choice/}{Education Next}). At the same time, child-rights and feminist legal scholars have raised sharp critiques of \textit{Pierce}'s parent-centered framework. In a landmark article, Barbara Bennett Woodhouse argued that \textit{Meyer} and \textit{Pierce} embedded a patriarchal conception of ``the child as property,'' treating children as vessels for parental ``destiny'' rather than independent rights-bearing persons with their own interests in an open future; more recent scholars such as James G. Dwyer have similarly questioned whether constitutionalizing robust parental control adequately protects children's independent educational and developmental welfare (\href{https://scholarship.law.gwu.edu/cgi/viewcontent.cgi?params=/context/faculty_publications/article/2970/&path_info=202401_Text.pdf}{GWU Law, citing Woodhouse, 33 Wm. \& Mary L. Rev. 995 (1992), and Dwyer (2025)}).

\textbf{Underlying Competing Values:} \textit{Pierce} negotiates a perennial three-way tension among (1) \textit{parental autonomy and cultural/religious pluralism}---shielding families and minority communities from majoritarian ideological conformity; (2) \textit{democratic statecraft}---the state's legitimate interest in forging a common civic culture and ensuring every child receives an adequate education; and (3) \textit{children's independent interests}---recognizing that children are neither ``mere creatures of the State'' nor mere extensions of their parents.

\section{Key Quotations}
\begin{quote}
``No question is raised concerning the power of the State reasonably to regulate all schools, to inspect, supervise and examine them, their teachers and pupils; to require that all children of proper age attend some school, that teachers shall be of good moral character and patriotic disposition, that certain studies plainly essential to good citizenship must be taught, and that nothing be taught which is manifestly inimical to the public welfare.''
\end{quote}
\noindent --- Justice McReynolds, Opinion of the Court, 268 U.S. at 534 (\href{https://supreme.justia.com/cases/federal/us/268/510/}{Justia}; \href{https://www.courtlistener.com/opinion/100694/pierce-v-society-of-sisters/}{CourtListener}).

\begin{quote}
``Under the doctrine of \textit{Meyer v. Nebraska}, 262 U.\,S. 390, we think it entirely plain that the Act of 1922 unreasonably interferes with the liberty of parents and guardians to direct the upbringing and education of children under their control. As often heretofore pointed out, rights guaranteed by the Constitution may not be abridged by legislation which has no reasonable relation to some purpose within the competency of the State.''
\end{quote}
\noindent --- Justice McReynolds, Opinion of the Court, 268 U.S. at 534--35 (\href{https://supreme.justia.com/cases/federal/us/268/510/}{Justia}).

\begin{quote}
``The fundamental theory of liberty upon which all governments in this Union repose excludes any general power of the State to standardize its children by forcing them to accept instruction from public teachers only. The child is not the mere creature of the State; those who nurture him and direct his destiny have the right, coupled with the high duty, to recognize and prepare him for additional obligations.''
\end{quote}
\noindent --- Justice McReynolds, Opinion of the Court, 268 U.S. at 535 (\href{https://supreme.justia.com/cases/federal/us/268/510/}{Justia}; \href{https://constitutioncenter.org/the-constitution/supreme-court-case-library/pierce-v-society-of-sisters}{National Constitution Center}).

\begin{quote}
``Plaintiffs asked protection against arbitrary, unreasonable and unlawful interference with their patrons and the consequent destruction of their business and property. Their interest is clear and immediate\dots''
\end{quote}
\noindent --- Justice McReynolds, Opinion of the Court, 268 U.S. at 536 (\href{https://supreme.justia.com/cases/federal/us/268/510/}{Justia}).

\end{document}
